% tk-02：辅助线速查目录——5 类代表辅助线一图速记
% A 三角形（中线倍长） B 梯形（平移一腰） C 圆（连半径+垂径）
% D 变换（旋转手拉手） E 坐标（过点作轴垂线）
\documentclass[tikz,border=4pt]{standalone}
\usepackage{ctex}
\usepackage{amsmath}
\usetikzlibrary{calc, decorations.markings}
\begin{document}
\begin{tikzpicture}[scale=0.8, thick,
  tick/.style={postaction={decorate, decoration={markings,
    mark=at position 0.5 with {\draw (-1.5pt,-3pt) -- (1.5pt,3pt);}}}}]

  % A. 三角形：中线倍长
  \begin{scope}[xshift=0cm, yshift=0cm]
    \coordinate (A) at (0.6,2.2);
    \coordinate (B) at (-1.4,0);
    \coordinate (C) at (1.8,0);
    \coordinate (D) at ($(B)!0.5!(C)$);
    \coordinate (E) at ($(A)!2!(D)$);
    \draw (A) -- (B) -- (C) -- cycle;
    \draw[red, dashed] (A) -- (E);
    \draw[red] (E) -- (C);
    \node[above] at (A) {\footnotesize $A$};
    \node[left]  at (B) {\footnotesize $B$};
    \node[right] at (C) {\footnotesize $C$};
    \node[below] at (D) {\footnotesize $D$};
    \node[below right] at (E) {\footnotesize $E$};
    \node[below] at (0.2,-1.0) {\small A. 三角形：中线倍长};
  \end{scope}

  % B. 梯形：平移一腰
  \begin{scope}[xshift=6cm, yshift=0cm]
    \coordinate (A) at (0.5,1.7);
    \coordinate (D) at (2.5,1.7);
    \coordinate (B) at (-0.4,0);
    \coordinate (C) at (3.2,0);
    \draw (A) -- (B) -- (C) -- (D) -- cycle;
    % 过 D 作 DE ∥ AB 交 BC 于 E
    \coordinate (E) at ($(B)+(D)-(A)$);
    \draw[red, dashed] (D) -- (E);
    \node[above left]  at (A) {\footnotesize $A$};
    \node[above right] at (D) {\footnotesize $D$};
    \node[below left]  at (B) {\footnotesize $B$};
    \node[below right] at (C) {\footnotesize $C$};
    \node[below] at (E) {\footnotesize $E$};
    \node[below] at (1.3,-1.0) {\small B. 梯形：平移一腰};
  \end{scope}

  % C. 圆：连半径 + 垂径
  \begin{scope}[xshift=12cm, yshift=0cm]
    \coordinate (O) at (1.3,0.8);
    \draw (O) circle (1.3);
    \coordinate (P) at ($(O)+({1.3*cos(160)},{1.3*sin(160)})$);
    \coordinate (Q) at ($(O)+({1.3*cos(20)},{1.3*sin(20)})$);
    \coordinate (M) at ($(P)!0.5!(Q)$);
    \draw (P) -- (Q);
    \draw[red, dashed] (O) -- (P);
    \draw[red, dashed] (O) -- (Q);
    \draw[red] (O) -- (M);
    \fill (O) circle (1.2pt);
    \node[above right] at (O) {\footnotesize $O$};
    \node[above left]  at (P) {\footnotesize $A$};
    \node[above right] at (Q) {\footnotesize $B$};
    \node[below] at (1.3,-0.95) {\small C. 圆：半径+垂径};
  \end{scope}

  % D. 变换：旋转构造（手拉手缩略）
  \begin{scope}[xshift=0cm, yshift=-5cm]
    \coordinate (O) at (1,0);
    \coordinate (A) at ({1+1.5*cos(150)},{1.5*sin(150)});
    \coordinate (B) at ({1+1.5*cos(110)},{1.5*sin(110)});
    \coordinate (C) at ({1+1.2*cos(40)},{1.2*sin(40)});
    \coordinate (D) at ({1+1.2*cos(0)},{1.2*sin(0)});
    \draw (O) -- (A) -- (B) -- cycle;
    \draw (O) -- (C) -- (D) -- cycle;
    \draw[red, dashed] (A) -- (C);
    \draw[red, dashed] (B) -- (D);
    \node[below] at (O) {\footnotesize $O$};
    \node[above left] at (A) {\footnotesize $A$};
    \node[above] at (B) {\footnotesize $B$};
    \node[above right] at (C) {\footnotesize $C$};
    \node[right] at (D) {\footnotesize $D$};
    \node[below] at (1,-1.0) {\small D. 变换：旋转手拉手};
  \end{scope}

  % E. 坐标：过点作轴垂线
  \begin{scope}[xshift=6cm, yshift=-5cm]
    \draw[->] (-0.4,0) -- (3.4,0) node[right] {\footnotesize $x$};
    \draw[->] (0,-0.4) -- (0,2.6) node[above] {\footnotesize $y$};
    \coordinate (P) at (2.4,1.8);
    \coordinate (Q) at (0.6,0.5);
    \fill (P) circle (1.5pt);
    \fill (Q) circle (1.5pt);
    \draw (Q) -- (P);
    \draw[red, dashed] (P) -- (2.4,0.5);
    \draw[red, dashed] (Q) -- (2.4,0.5);
    \node[above right] at (P) {\footnotesize $P$};
    \node[above left]  at (Q) {\footnotesize $Q$};
    \node[below] at (1.5,-1.0) {\small E. 坐标：作轴垂线};
  \end{scope}

\end{tikzpicture}
\end{document}
